\chapter{Mind Your Motive}

These notes follow Michael Sandel's presentation of Kant as the point in the course where morality and freedom must be thought together. The lecture proceeds in a deliberate order. We begin with Kant's challenge to utilitarianism, move from there to a demanding idea of freedom, then ask what gives an action moral worth, and only after those preparations do we reach the categorical imperative itself. The chapter keeps that order, preserves the main classroom objections, and retains the key visual evidence. It is prepared for the LazyLearn track, with curation by LazyingArt LLC.

\section{Why Kant, and Why Now}

Sandel opens by warning us that Kant is the hardest philosopher in the course. The warning is not theatrical. It is meant to prepare us for a style of moral argument that does not begin from happiness, preference, or ownership, but from reason and freedom. He also briefly humanizes Kant before beginning the argument proper: the young prodigy at K\"onigsberg, the underpaid lecturer, the late but extraordinary author of the \emph{Critique of Pure Reason} and the \emph{Groundwork for the Metaphysics of Morals}. The biographical sketch matters because it gives the audience a moment of orientation before the difficult turn.

The philosophical stake is then announced in two large questions. Kant's moral philosophy is said to concern
\begin{enumerate}
\item the supreme principle of morality, and
\item the possibility of freedom.
\end{enumerate}
This is already a clue about the architecture of the lecture. We are not being given one topic after another. We are being led through an argument in which the account of morality will have to match the account of freedom.

The first move is anti-utilitarian. Kant does not deny that we seek pleasure and avoid pain. In that limited sense, Sandel says, the utilitarians are half right. But they are wrong to claim that pleasure and pain are our sovereign masters. What distinguishes human beings is not merely that we can suffer or enjoy, but that we are rational beings capable of free action.

The lecture's first structural compression is therefore
\begin{equation}
\text{rational capacity} \Longrightarrow \text{dignity and respect}.
\end{equation}
This is why, for Kant, rights do not rest on self-ownership, and why respect for persons does not depend on an aggregate calculus of happiness. If rational agency is what gives human beings their distinctive worth, then persons may not be used merely as instruments, even for desirable ends.

\section{Rational Beings, Dignity, and Freedom Beyond Desire}

From dignity Sandel turns to freedom. He first states the common view in order to set it aside. Ordinarily we think of freedom as doing what we want, or at least as facing no obstacle to the satisfaction of our wants. Kant's point is that this account is too shallow, because it leaves the will under the governance of desires it did not choose.

The lecture's contrast is sharp:
\begin{align}
\text{freedom} &\neq \text{doing what we want},\\
\text{freedom} &= \text{autonomy}.
\end{align}
And autonomy immediately requires an opposite:
\begin{align}
\text{autonomy} &= \text{acting according to a law I give myself},\\
\text{heteronomy} &= \text{acting according to an inclination or desire not self-given}.
\end{align}

Sandel then motivates the distinction with examples rather than definitions alone. The Sprite slogan, ``Obey your thirst,'' becomes a compressed lesson in Kantian moral psychology. If thirst governs the action, then the source of the command is not reason but appetite. The action may feel voluntary, but it is not free in the demanding sense Kant cares about. We are not yet authors of the law of our action.

In the same spirit, the billiard-ball example clarifies what Kant means by necessity. A billiard ball falls according to gravity; we do not say it acts freely. The force of the example is that when appetite and impulse determine our will in an equally law-governed way, the will is still being pushed along by nature. It is not yet legislating for itself.

The lecture's causal map can therefore be written as
\begin{equation}
\text{inclination} \to \text{heteronomy} \to \text{necessity},
\qquad
\text{reason} \to \text{autonomy} \to \text{freedom}.
\end{equation}
And the immediate moral consequence is already visible:
\begin{equation}
\text{act on appetite or desire} \Longrightarrow \text{natural necessity, not freedom}.
\end{equation}
Once this is in view, Sandel can say that autonomous action is action for its own sake rather than action as an instrument of externally given ends:
\begin{equation}
\text{autonomous action} \Longrightarrow \text{act for its own sake, as an end in itself}.
\end{equation}

\subsection{Question \& Answer}

\paragraph{Question.} If satisfying my desires feels free, why does Kant say it is a form of unfreedom?

\paragraph{Answer.} Because the feeling of spontaneity is not yet self-legislation. I may choose this drink rather than that one, or this meal rather than that one, but if what governs the choice is an appetite I simply find in myself, then the governing principle is not mine in the relevant sense. Kant's standard is stricter: freedom is not efficient desire-satisfaction but authorship of the law according to which we act.

This is also why Sandel can move from freedom back to dignity without changing the topic. If we are capable of legislating for ourselves through reason, then we are not merely loci of appetite. We are beings who can stand back from appetite, evaluate it, and act from another source.

\section{Moral Worth, Good Will, and the Motive of Duty}

Once freedom has been redefined, morality must be redefined with it. Sandel makes the transition explicitly: if autonomy means self-legislation, then the moral worth of an action cannot lie in consequences alone. It must lie in the will from which the action proceeds.

The basic claim is
\begin{equation}
\text{moral worth} \Longrightarrow \text{depends on motive, not consequences}.
\end{equation}
What matters is not simply that the action conforms to what morality requires, but that it is done \emph{for the sake of} what morality requires. This is the point of Kant's account of the good will:
\begin{equation}
\text{good will} \Longrightarrow \text{good in itself}.
\end{equation}
A will can fail to accomplish what it aims at and still retain its full moral worth, because its worth is not borrowed from outcomes.

Sandel then introduces the lecture's central moral test:
\begin{equation}
\text{right action} + \text{duty as motive} \Longrightarrow \text{moral worth}.
\end{equation}
And, conversely,
\begin{equation}
\text{right action} + \text{self-interest alone} \not\Longrightarrow \text{moral worth}.
\end{equation}
This is also the point at which the lecture's vocabulary settles into a stable opposition:
\begin{equation}
\text{duty} \leftrightarrow \text{inclination}.
\end{equation}
Inclination here is a broad category. It includes appetite, desire, sympathy, altruism, prudence, self-interest, and any other contingent motive that comes to us as given.

This gives the lecture its next transition. We have moved from a theory of freedom to a theory of the will. The question is no longer merely, ``What did the action produce?'' It is now, ``From what principle was the action done?''

\section{Examples That Isolate Motive}

At this point Sandel does not continue with abstractions. He slows down and tests the claim with cases. The sequence matters, because each example holds the outward act fixed while varying the motive.

\paragraph{The shopkeeper.} The shopkeeper gives correct change to an inexperienced customer. Outwardly he does the right thing. But his motive is prudential: cheating might damage his reputation and cost him business. The action conforms to duty; it is not done from duty. This is the lecture's cleanest example of the gap between external correctness and moral worth.

\paragraph{The duty against suicide.} Here Kant's problem is that ordinary life mixes motives too freely. Most people who continue living have many reasons for doing so. To isolate the relevant motive, Kant imagines a miserable person who nevertheless preserves himself because he recognizes a duty to do so. Again the moral point is not what happens but what explains the action.

\paragraph{The Better Business Bureau and the honor-code discount.} These examples widen the frame without changing the logic. ``Honesty is the best policy'' is not yet a moral principle in Kant's sense if what makes honesty attractive is profitability. Likewise, an honor pledge sustained by commercial discounts simply reproduces the shopkeeper case in another institutional setting.

Sandel summarizes the connection between morality and freedom by saying that only when we act from duty, resisting inclination and self-interest, do we act freely in the strong Kantian sense. In the lecture's compressed form,
\begin{equation}
\text{act from duty} \Longrightarrow \text{autonomy},
\qquad
\text{act from inclination} \Longrightarrow \text{heteronomy}.
\end{equation}

\subsection{Question \& Answer}

\paragraph{Question.} Can an action still have moral worth if duty is accompanied by self-interest, sentiment, or the desire to be moral?

\paragraph{Answer.} Sandel's answer is careful. If self-interest is the ground of the action, moral worth is lost. But the mere presence of feeling does not by itself destroy moral worth. Kant allows that there must be some incentive to obey the moral law, yet it cannot be a self-interested incentive on pain of defeating the whole account. This is why Sandel introduces reverence for the moral law at precisely this moment. Feelings may accompany an act, and may even support it, but they are not what finally justifies it. In compact form,
\begin{equation}
\text{supporting sentiments} \neq \text{ground of moral worth}.
\end{equation}
This is the refinement that prevents the view from becoming a caricature. The issue is not the presence or absence of sentiment; the issue is which motive gives the action its moral worth.

\section{One Moral Law or Many? Shared Reason and the Internal Reset}

The next classroom objection goes deeper. If autonomy means acting according to a law I give myself, why does morality not collapse into subjectivity? Why are there not as many moral laws as there are persons?

Sandel treats this as a real problem rather than a misunderstanding. If the moral law were merely private conscience, Kantian autonomy would indeed fragment into idiosyncrasy. The answer is that the source of the law is not my empirical personality with its accidental history and preferences. The source is the rational capacity we share as rational beings.

The lecture's answer is therefore
\begin{equation}
\text{shared practical reason} \Longrightarrow \text{one moral law, not many private laws}.
\end{equation}
This is why Sandel insists that practical reason is not biographical. It is not the set of reasons that arise from my upbringing, interests, or temperament. It is the universal rational capacity that exists, as he says, undifferentiated in all of us.

\subsection{Question \& Answer}

\paragraph{Question.} If the moral law is self-given, why is morality not merely subjective?

\paragraph{Answer.} Because the relevant self is not the contingent self. Kantian self-legislation is not ``I invent my own morality.'' It is the rational will giving itself the law that any rational being, insofar as it is rational, must be able to acknowledge. Sandel therefore moves from autonomy to universality, not from autonomy to relativism.

At this point the lecture pauses and then deliberately resets. Sandel returns to Kant, restates the two big questions---the supreme principle of morality and the possibility of freedom---and tells us that one way through the \emph{Groundwork} is to keep a set of linked contrasts constantly in view. This reset matters. It gathers the earlier material into a tighter structure and prepares the final movement of the lecture.

\section{Three Contrasts and the Imperative Distinction}

Sandel now reorganizes the lecture around three dualisms:
\begin{equation}
\text{duty} \leftrightarrow \text{inclination},
\qquad
\text{autonomy} \leftrightarrow \text{heteronomy},
\qquad
\text{categorical imperative} \leftrightarrow \text{hypothetical imperative}.
\end{equation}
He briefly revisits the first contrast with the spelling-bee anecdote. The point of the story is not the newspaper detail but the problem of mixed motive. The student's confession appears admirable, but it is morally complicated by his admission that he did not want to feel like a slime. Judith's defense of Kant is exactly the right one for Sandel's purposes: more than one motive may be present, but moral worth depends on whether duty supplies the reason for acting.

The second contrast then returns: the will can be determined autonomously or heteronomously. Only in the first case is it free. And this leads straight to the third contrast, the contrast between the two kinds of imperative reason can issue.

A hypothetical imperative has the instrumental form
\begin{equation}
\text{If you want } X,\ \text{do } Y.
\end{equation}
The lecture's example is straightforward:
\begin{equation}
\text{If you want a good business reputation, do not shortchange your customers.}
\end{equation}
This is a command of reason, but only conditional on a further end.

A categorical imperative, by contrast, commands without dependence on any further purpose. The lecture's slide captures this contrast especially well and is worth preserving as evidence of the lecture's structure.

\begin{figure}[htbp]
\centering
\includegraphics[width=0.82\textwidth]{lecture_06_figure_02.png}
\caption{Kant on hypothetical and categorical imperatives.}
\end{figure}

The safest reconstruction, faithful to the visible wording while avoiding the omitted portions of the quotation, is
\begin{align}
\text{good merely as means} &\Longrightarrow \text{hypothetical imperative},\\
\text{good in itself} &\Longrightarrow \text{categorical imperative}.
\end{align}
This is the point at which the lecture's pieces lock together. If morality requires duty rather than inclination, and if freedom requires autonomy rather than heteronomy, then the command appropriate to moral freedom must be categorical rather than hypothetical.

\section{Universal Law, Humanity as End, and the Reach of Respect}

Only now does Sandel ask what the moral law actually commands. He notes that Kant offers several formulations of the categorical imperative, but the lecture concentrates on two.

The first is the formula of universal law. A maxim is defined as the rule or principle explaining why we are acting. The command is:
\begin{equation}
\begin{aligned}
\text{Act only on that maxim whereby you can at the same time}\\
\text{will that it should become a universal law.}
\end{aligned}
\end{equation}

\paragraph{Worked example: false promise.} Sandel's worked case is the false promise. We need money, we know we cannot repay promptly, and we nevertheless promise repayment in order to secure a loan. The relevant maxim may be written schematically as
\begin{equation}
\text{When I need money, I will promise repayment though I know I cannot repay.}
\end{equation}
The argument then unfolds step by step:
\begin{enumerate}
\item Formulate the maxim that explains the intended action.
\item Universalize it: suppose everyone acted on that maxim whenever in need.
\item Ask whether the institution presupposed by the maxim could survive that universalization.
\item In this case, universalized false promising destroys the practice of promising itself.
\item The maxim therefore defeats itself and cannot be willed as a universal law.
\end{enumerate}
The lecture's contradiction structure is
\begin{equation}
\text{false-promise maxim universalized} \Longrightarrow \text{promise practice collapses} \Longrightarrow \text{contradiction}.
\end{equation}

\subsection{Question \& Answer}

\paragraph{Question.} Does this test secretly rely on consequences after all?

\paragraph{Answer.} Sandel pauses because the objection is natural. It can sound as if Kant is saying, ``Do not lie because the world will go worse if everyone lies.'' But that would be consequentialist reasoning, and Sandel insists that this is not the point. Universalization is a test, not the moral reason itself. The point is to see whether my maxim can be justified without privileging my special interests over everyone else's, and whether it can stand as a law for rational agents as such.

Sandel then turns to the second formulation because it brings the moral intuition into sharper focus. A categorical imperative cannot be grounded in merely private or contingent ends. Its ground must lie in something with absolute value. Kant's answer is that rational beings exist as ends in themselves. The formula is:
\begin{equation}
\begin{aligned}
\text{Treat humanity, whether in yourself or another, never simply as a means}\\
\text{but always at the same time as an end.}
\end{aligned}
\end{equation}

The false promise now becomes immediately transparent. In deceiving another for the sake of money, I am not merely interacting with a person in pursuit of a purpose. I am manipulating his rational agency and subordinating it to my own private end.

The same structure is then applied to murder and suicide:
\begin{equation}
\text{murder},\ \text{suicide} \Longrightarrow \text{use humanity merely as means} \Longrightarrow \text{violate dignity}.
\end{equation}
This is why Sandel says that, for Kant, suicide is morally on a par with murder. In both cases what is violated is humanity itself, the rational nature that commands respect. Because that rational nature is universal rather than particular, Kantian respect differs from love, sympathy, solidarity, or altruism. Those may attach to particular people because of who they are. Respect, by contrast, is respect for rational being as such.

\subsection{Question \& Answer}

\paragraph{Question.} How can we use other people as means in everyday life without treating them merely as means?

\paragraph{Answer.} Patrick's objection is exactly the right one to raise here, because ordinary life is full of purposive interaction. We buy, sell, hire, ask directions, accept help, and cooperate in projects. Kant's prohibition is not against all such use. The crucial distinction is
\begin{equation}
\text{use others as means} \neq \text{use others merely as means}.
\end{equation}
Ordinary social and market interaction is compatible with the categorical imperative when it proceeds under conditions of respect for the other as a rational being. Fraud, coercion, and manipulation are wrong because they erase that distinction. The other person ceases to be a participant in the action and becomes only raw material for my purposes.

\section{Summary}

Sandel's lecture has a strict sequence, and the sequence is philosophically necessary. We begin with Kant as the critic of utilitarianism, then discover that the ground of dignity is rational agency. From there freedom is redefined as autonomy rather than desire-satisfaction. Once freedom is understood in that way, morality can no longer be based on consequences alone; it must concern motive, and specifically the motive of duty. The contrasts duty versus inclination, autonomy versus heteronomy, and categorical versus hypothetical imperative then emerge as three faces of the same structure.

The two central formulations of the categorical imperative complete the argument. The formula of universal law tests whether a maxim can be willed universally; the formula of humanity as an end states more directly what the moral law requires of us. By the end of the lecture we do not simply have a list of duties. We have an entire moral architecture: freedom as self-legislation, moral worth as action from duty, and respect for persons as the refusal to treat rational beings merely as means. The question Sandel leaves open is the right one to carry forward: not whether Kant's morality is easy, but whether its severity may be the price of taking freedom and dignity seriously.